\section{Varying projective readout and conditional preparations}\label{sec:readout72}
The fixed input basis can be allowed to vary when its outcome remains
observable. This changes the coding scale: destruction of the measured
quantum input does not make the readout coordinate a preparation
coordinate. The distinguishing experiment may use entangled inputs.

Let $\mathfrak F_d$ be the space of ordered orthogonal rank-one
projections $P=(P_1,\ldots,P_d)$ on $\C^d$ with $\sum_xP_x=I_d$.
It is the flag manifold $U(d)/\mathbb T^d$, of real dimension
$b=d^2-d$. The column phases are immaterial; permutation of the
ordered outcomes is not a gauge. For each $x$ let
$\sigma_x=(\sigma_{xy})_{y=1}^m\in\mathcal S_{m,n}$ and set
\begin{equation}\label{eq:readoutfamily72}
 \mathcal F_{P,\sigma,(x,y)}(X)=\tr(P_xX)\sigma_{xy},\qquad
 J_{xy}=P_x^{\mathsf T}\otimes\sigma_{xy}.
\end{equation}
Both labels $(x,y)$ are retained. The transpose in this input-first
Choi formula is essential, including for complex projections. Fix
public bounds $0\le r_{xy}\le n$, with at least one positive bound
in each row, and put
\[
 v_x=\sum_y r_{xy}(2n-r_{xy})-1,\qquad V=\sum_xv_x.
\]
The family $\mathfrak R_{\boldsymbol r}$ consists of
\eqref{eq:readoutfamily72} with these rank bounds. Let
$\mathcal R_N(\delta)$ be its covering number for $d_N$, allowing
arbitrary legal memoryless instrument centres of the same interface;
$\mathcal R_N^{\Q}(\delta)$ requires rational centres. Public dimensions,
outcome labels, ranks, horizon and tolerance are not payload.

\begin{theorem}[Mixed-scale coding with observable readout]\label{thm:readoutentropy72}
For fixed $d,n,m,\boldsymbol r$ there are constants $c,C>0$ such that
for $N\ge1$ and $0<\delta\le1/32$,
\begin{equation}\label{eq:readoutcover72}
 c(N/\delta)^b(\sqrt N/\delta)^V
 \le\mathcal R_N(\delta)
 \le\mathcal R_N^{\Q}(\delta)
 \le C(N/\delta)^b(\sqrt N/\delta)^V.
\end{equation}
Consequently the optimal reusable description length is
\begin{equation}\label{eq:readoutbits72}
 (b+V/2)\log_2N+(b+V)\log_2(1/\delta)+O(1),
\end{equation}
with remainder uniform in $N$ and $\delta$. The upper centres belong
to the displayed family, preserve its zero conditional blocks and do
not increase any conditional rank. For supplied rational projections
and rational row states there is a finite exact rational encoder,
without choosing algebraic unit-vector representatives or searching
an exhaustive unitary net. Real targets admit the same rational-centre
cover by the exact-real construction. If $b+V=0$, the family is a
known singleton and needs zero payload.
\end{theorem}
The error range here is not the full submaximal range of
Theorem~\ref{thm:cqentropy71}. The readout part of the lower proof uses
pairwise separation; it is not an application of the full-error
half-dimensional estimator theorem.

\subsection{The two distinguishing scales}
Write $\mathsf M_P(X)=(\tr(P_xX))_x$ and
$q(P,Q)^2=\sum_x\|P_x-Q_x\|_F^2$.

\begin{lemma}[Readout modulus]\label{lem:readoutmodulus72}
If $z=d_1(\mathsf M_P,\mathsf M_Q)$, then
\begin{equation}\label{eq:readoutone72}
 \sqrt{2/d}\,q(P,Q)\le z\le\sqrt d\,q(P,Q).
\end{equation}
For $z<2$ one has
\begin{equation}\label{eq:readoutmany72}
 d_N(\mathsf M_P,\mathsf M_Q)
   =2\sin\!\left(\min\{N\arcsin(z/2),\pi/2\}\right).
\end{equation}
For $z=2$ the distance is two at every $N\ge1$. In particular it is
at least $(2/\pi)\min\{Nz,1\}$ and at most $Nz$.
\end{lemma}
\begin{proof}
The difference of two rank-one projections has nonzero eigenvalues
$\pm s$, so its Frobenius norm is $\sqrt2s$. An input eigenvector
maximizing one event difference gives total unhalved outcome distance
at least $2\max_x\|P_x-Q_x\|_\infty$. This proves the lower bound.
For an arbitrary reference input state $\rho$, duality gives
\[
 \big\|\tr_A(((P_x-Q_x)\otimes I)\rho)\big\|_1
 \le\|P_x-Q_x\|_\infty.
\]
Indeed a Hermitian contraction $H$ on the reference tests this block
by $\tr(((P_x-Q_x)\otimes H)\rho)$. Sum over the classical blocks
and apply Cauchy--Schwarz to obtain the upper bound in
\eqref{eq:readoutone72}.

We use here the exact measurement-discrimination theorem of
Pucha\l a--Pawela--Krawiec--Kukulski--Oszmaniec
\cite[Corollary~1 and Theorem~2]{PPKKO72}. If $U$ represents the relative
ordered basis, their optimized eigenvalue-arc length $\Upsilon(U)$
gives distance $2\sin(N\Upsilon(U)/2)$ until
$N\Upsilon(U)\ge\pi$, when the distance is two; their Theorem~2
allows general adaptive networks. At one use, $z<2$ implies
$\Upsilon(U)=2\arcsin(z/2)$. The already-perfect case is treated
separately and need not be inverted. This proves
\eqref{eq:readoutmany72}. A publicly stopped tester is contained in
a fixed $N$-slot network by performing irrelevant calls after the
stop and discarding their outputs. Conversely, parallel testers are
admissible here. The elementary bound
$\sin t\ge2t/\pi$ for $0\le t\le\pi/2$ gives the stated lower
modulus; the upper bound also follows directly by a hybrid argument.
The exact multiple-use theorem is an imported result, not an inference
from a finite-dimensional spectral gap.
\end{proof}

\begin{lemma}[Measurement and programme comparisons]\label{lem:readoutmetric72}
For arbitrary row states,
\[
 d_N(\mathcal F_{P,\sigma},\mathcal F_{Q,\tau})
          \ge d_N(\mathsf M_P,\mathsf M_Q).
\]
When $P=Q$, the distance is at least
$\max_x\|\Omega_x^{\otimes N}-\Lambda_x^{\otimes N}\|_1$,
where $\Omega_x=\bigoplus_y\sigma_{xy}$ and
$\Lambda_x=\bigoplus_y\tau_{xy}$. For purifications whose rowwise
vector differences have lengths $\eta_x$, one has
\begin{equation}\label{eq:readoutupper72}
 d_N(\mathcal F_{P,\sigma},\mathcal F_{Q,\tau})
 \le N d_1(\mathsf M_P,\mathsf M_Q)
                      +2\sqrt{N\sum_x\eta_x^2}.
\end{equation}
\end{lemma}
\begin{proof}
Discarding the quantum outputs and $y$ gives the measurement channel,
so every measurement tester is a tester of the displayed instruments.
For a common flag, feeding a unit vector in $\operatorname{ran}P_x$
at every call gives the $x$th product-state witness. For the upper
bound first change $P$ to $Q$ while keeping $\sigma$ fixed. Appending
the conditional row state is a common postprocessing of the readout,
and a slotwise hybrid gives the first term. Then change the row states
with $Q$ fixed. At each slot a common processor measures $Q$, selects
one of the supplied row programme states, retains $x$, and discards
the unused rows. The product-purification argument of
Lemma~\ref{lem:cqmetric71} gives the second term, also with references,
feedback and stopping. This is the standard common-environment method
\cite[Proposition~2]{DW72}; retaining the readout label changes neither
its positivity nor its contraction argument. The upper comparison is
not asserted to equal either lower comparison.
\end{proof}

\subsection{A bounded rational atlas without column phases}
There is a direct flag codec, rather than an exhaustive search over
unitaries. We give its quantitative bound because it also controls
singular rational descriptions whose unit vectors are not rational.

\begin{lemma}[Pivoted triangular flag charts]\label{lem:flagatlas72}
Every $P\in\mathfrak F_d$ is obtained from an ordered basis of columns
of $\Pi^{-1}L$, where $\Pi$ is a permutation matrix and $L$ is unit
lower triangular with complex entries of modulus at most one below
the diagonal. Its ordered rank-one projections are the successive
differences of projections onto the first $j$ columns. Thus $d!$
charts, each with $b=d^2-d$ real coordinates, suffice. Given rational
$P_x$, both the chart and its coordinates can be found rationally.

Round each real and imaginary coordinate of $L$ towards zero at grid
$B\ge1$ and let $\widehat P$ be the decoded flag. Put
\[
 K_d=(2d)^{d-1},\qquad C_d=4d^2K_d.
\]
Then every decoded flag is exactly legal and rational, and
\begin{equation}\label{eq:flaground72}
 d_N(\mathsf M_P,\mathsf M_{\widehat P})\le C_dN/B.
\end{equation}
There are at most $d!(2B+1)^b$ codewords. The expanded projection
entries have $O_d(1+\log B)$ bits.
\end{lemma}
\begin{proof}
Choose a nonzero column $v_x=P_xe_{k_x}$, taking the first positive
diagonal position. These columns are rational whenever the input
projections are rational. The columns
are nonzero and mutually orthogonal; hence $A=[v_1\ \cdots\ v_d]$
is invertible. Gaussian elimination with partial row pivoting gives
$\Pi A=LR$, with $L$ unit lower triangular and $R$ upper triangular
and invertible. At every pivot choose a residual entry of largest
modulus, breaking ties by row order. Each multiplier below it has
modulus at most one. Subsequent row interchanges interchange the
already recorded multipliers as well. All comparisons of squared
moduli and all arithmetic are rational for rational $A$.
Right multiplication by an invertible upper triangular matrix
preserves each initial column span. Thus the first $j$ columns of
$\Pi^{-1}L$ span $\operatorname{ran}(P_1+\cdots+P_j)$.
This proves the chart assertion without choosing square roots.

Any unit lower triangular matrix in the coordinate cube
$|\Re L_{ij}|,|\Im L_{ij}|\le1$ is invertible, even if it was not
produced by elimination. If $Z_j$ is its first $j$ columns, put
\[
 E_j=Z_j(Z_j^*Z_j)^{-1}Z_j^*,\qquad E_0=0.
\]
The differences $E_j-E_{j-1}$ are mutually orthogonal rank-one
projections with sum $I$. Permute their coordinates back by $\Pi$.
The same formulas decode the rounded rational matrix exactly.

Throughout the cube, $\|L\|_2\le2d$ and $\det L=1$, so its least
singular value is at least $K_d^{-1}$. Restriction to any first $j$
columns preserves this lower bound. For two such matrices $L,L'$,
write $E_j,E'_j$ for their column-space projections. Using
$(I-E'_j)Z'_j=0$ and $E_j=Z_jZ_j^\dagger$ gives
\[
 \|(I-E'_j)E_j\|_F\le K_d\|Z_j-Z'_j\|_F.
\]
The analogous bound with the primes interchanged and the identity
$E_j-E'_j=(I-E'_j)E_j-E'_j(I-E_j)$ imply
$\|E_j-E'_j\|_F\le2K_d\|L-L'\|_F$.
Consequently each rank-one difference is at most
$4K_d\|L-L'\|_F$. The one-use block estimate used above bounds the
measurement distance by the sum of these norms. Rounding changes
$L$ by at most $\sqrt b/B\le d/B$, proving the one-use bound
$C_d/B$ and then \eqref{eq:flaground72} by hybridization.

Record the permutation and the $b$ signed digits in $[-B,B]$.
No phases or rational unit-vector normalizations are encoded. All
inverse Gram matrices in the displayed decoding formula have fixed
size and rational entries of degree bounded in terms of $d$ in the
digits and $B$. The determinant is nonzero. Clearing denominators
and the adjugate formula bound every numerator and denominator by
$B^{O_d(1)}$ times a fixed constant. This proves the bit-length claim.
\end{proof}
Partial pivoting is used here only to select a bounded chart.
The decoder itself needs no pivot test on the target and always
returns a flag for any valid chart word. A target-bound verifier
nevertheless checks the entire canonical encoding, not just legality.

\subsection{Proof of the covering law}
\begin{proof}[Proof of Theorem~\ref{thm:readoutentropy72}]
For the upper bound, use the flag code at
$B_f=\max\{1,\lceil2C_dN/\delta\rceil\}$ and the conditional factor
code at $B_s=\max\{1,\lceil8\sqrt{NV}/\delta\rceil\}$.
The two terms of \eqref{eq:readoutupper72} are each at most
$\delta/2$. The factor code counts every actual pivot mask, anchor
and sign. Its total cardinality is at most $C(2B_s+1)^V$.
Multiply by $d!(2B_f+1)^b$. On the stated range this proves the
upper bound in \eqref{eq:readoutcover72}. Each factor row is trace
normalized and rank nonincreasing. Formula \eqref{eq:readoutfamily72}
then gives positive Choi blocks with joint partial trace $I_d$;
their ranks equal those of the conditional states. Rationality and
zero preservation follow exactly. Real target projectors admit the
same chart selection and rounding; only the algorithmic assertion
about rational targets uses finite rational comparisons.

For the lower bound, choose a small closed coordinate box around a
flag in $\mathfrak F_d$ on which $q$ is bi-Lipschitz to a Euclidean
metric in $b$ dimensions. Such a chart can also be checked directly
from $L=I+T$: to first order the off-diagonal entries of the ordered
projections recover every real and imaginary entry of strictly lower
triangular $T$. The derivative is injective, and shrinking to a compact
box gives fixed forward and inverse Lipschitz constants. Grid-pack it
at $q$-separation greater than $A_d\delta/N$, where $A_d$ is fixed
large enough that Lemma~\ref{lem:readoutmodulus72} gives distance
strictly greater than $2\delta$. Indeed its constant lower plateau
$2/\pi$ exceeds $2\delta$ on this range. The packing has at least
$c(N/\delta)^b$ points after reducing $c$ to include the coarse-grid
cases.

In every nontrivial row choose the compact maximal-rank factor patch
from Theorem~\ref{thm:rankentropy68}. The product row patch has
$V$ real dimensions. Pack it in tuple Frobenius distance at more
than $A_{d,n,m,\boldsymbol r}\delta/\sqrt N$. Some row of two
distinct grid points then differs by at least the tuple distance
 divided by $\sqrt d$. The local product-Choi test
(Corollary~\ref{cor:localchoi68}, with input dimension one) implies
that its $N$-copy states are separated by more than $2\delta$.
The row packing has cardinality at least $c(\sqrt N/\delta)^V$.

Take the Cartesian product of these two finite packings. If two
points have different flags, discarding the conditional outputs gives
the separating measurement witness regardless of their rows. If their
flags agree, the product-state witness in a differing row separates
them. These are legitimate pair-dependent testers, not incompatible
tests assumed to hold in one common experiment. Hence all distinct
product points have $d_N$-distance greater than $2\delta$. By the
triangle inequality, a ball centred at any legal memoryless instrument
covers at most one of them. Multiplying cardinalities proves the
lower bound. A fixed-dimensional ceiling in the word length changes
only the additive constant. Finally, $b+V=0$ forces $d=1$ and a
singleton rank-state family, proving the last assertion.
\end{proof}

\begin{corollary}[Readout recovered from orthogonal output supports]\label{cor:observablereadout72}
Suppose the label $x$ in \eqref{eq:readoutfamily72} is not output,
but the $x$th row is supported in a fixed public subspace $H_x$,
with the $H_x$ mutually orthogonal. Then the same coding law holds
with $n$ replaced by $n_x=\dim H_x$ in each row dimension. For
public rational coordinate embeddings the upper code is rational.
\end{corollary}
\begin{proof}
Measuring the fixed output projections onto $H_x$ recovers $x$
without disturbing the target's conditional state. This operation
and forgetting $x$ are fixed CPTP postprocessings inverse on the
family. Inserting them at every slot proves equality of the two
metrics on target pairs, with references and adaptive feedback.
The preceding upper code and product packing therefore apply. The
lower packing still allows arbitrary legal centres. We do not infer
a target-independent input retraction for varying flags, or a rank
bound on the postprocessing of an arbitrary centre.
\end{proof}

\begin{remark}[What can disappear when readout is unobservable]\label{rem:unflagged72}
If the label $x$ is discarded and all conditional rows equal one
state tuple $\sigma$, then
$\sum_x\tr(P_xX)\sigma_y=\tr(X)\sigma_y$ for every flag.
Thus the entire basis coordinate disappears exactly. This example
separates an observable readout theorem from an unjustified parameter
count for arbitrary unlabelled overlapping output rows. The two
proved families above do have the asserted mixed scale; neither
statement classifies a general coupled coherent/support-changing
instrument boundary.
\end{remark}
